% ============================================================================
% CAPÍTULO 1: INTRODUCCIÓN
% ============================================================================
% 
% INSTRUCCIONES GENERALES:
% - Este capítulo presenta una visión general de la investigación
% - Debe ser claro, conciso y motivar al lector
% - Extensión típica: 10-15 páginas
% - Reemplaza TODO el texto entre [corchetes]
% ============================================================================

\chapter{Introducción}
\label{cap:introduccion}

% ----------------------------------------------------------------------------
% CONTENIDO REQUERIDO PARA INTRODUCCIÓN:
% ----------------------------------------------------------------------------
% Este capítulo debe incluir una visión general de:
% - El contexto y relevancia del tema de investigación
% - La problemática que abordas
% - Una breve descripción de lo que harás
% - La estructura del documento
%
% NOTA IMPORTANTE: La Delimitación del Objeto de Estudio (planteamiento del
%                  problema, pregunta de investigación, hipótesis, objetivos)
%                  va en el CAPÍTULO 3, NO aquí.
% ----------------------------------------------------------------------------

[Escribe aquí tu introducción. Debe presentar el contexto general, la problemática y el alcance de tu investigación sin entrar en detalles de objetivos o hipótesis, ya que eso va en el Capítulo 3.]

\textbf{Sugerencia de estructura:}

\textbf{1. Contexto General (2-3 párrafos):}

Introduce el tema de tu investigación explicando por qué es relevante. Presenta el panorama general del área de estudio.

\textbf{2. Problemática (2-3 párrafos):}

Describe brevemente el problema o vacío que has identificado en la literatura o en la práctica profesional.

\textbf{3. Alcance de la Investigación (1-2 párrafos):}

Indica de manera general qué abordarás en tu tesis, qué tipo de metodología usarás (sin detalles) y qué aportación esperas hacer.

\textbf{4. Organización del Documento (1 párrafo):}

% ----------------------------------------------------------------------------
% SECCIÓN: Organización del Documento
% ----------------------------------------------------------------------------
\section*{Organización del Documento}

Este documento de tesis se organiza de la siguiente manera:

\begin{itemize}
    \item \textbf{\Cref{cap:marco_teorico}} presenta el marco teórico y antecedentes, incluyendo la revisión de literatura actualizada, fuentes primarias y análisis crítico de investigaciones previas relacionadas con el tema.
    
    \item \textbf{\Cref{cap:delimitacion}} establece la delimitación del objeto de estudio, presentando el planteamiento del problema, la pregunta de investigación, las hipótesis y los objetivos (general y específicos) de la investigación.
    
    \item \textbf{\Cref{cap:justificacion}} justifica la realización de esta investigación considerando su conveniencia, relevancia social, implicaciones prácticas, valor teórico y utilidad metodológica.
    
    \item \textbf{\Cref{cap:metodologia}} describe detalladamente los materiales y métodos empleados, incluyendo el diseño del estudio, la población, las variables, las técnicas, los procedimientos y las consideraciones éticas.
    
    \item \textbf{\Cref{cap:resultados}} presenta los resultados obtenidos, incluyendo la validación de instrumentos, caracterización de la población, análisis univariado, bivariado y multivariado, así como la tesis derivada de la investigación.
    
    \item \textbf{\Cref{cap:discusion}} realiza aportaciones al conocimiento a partir de los resultados más relevantes, discutiendo los hallazgos en relación con otras investigaciones, identificando logros, riesgos y limitaciones.
    
    \item \textbf{\Cref{cap:conclusiones}} describe de forma concreta los resultados que contribuyen al logro de los objetivos e hipótesis con sustento estadístico.
\end{itemize}

Finalmente, se presentan las referencias bibliográficas y los anexos con material complementario.

% ============================================================================
% NOTAS FINALES:
% ============================================================================
% - Recuerda citar fuentes relevantes usando \cite{clave_referencia}
% - Mantén un tono académico pero accesible
% - Evita entrar en detalles metodológicos (eso va en Cap. 5)
% - No incluyas objetivos ni hipótesis aquí (eso va en Cap. 3)
% ============================================================================
